\documentclass[11pt]{article}

\usepackage[T1]{fontenc}
\usepackage[utf8]{inputenc}
\usepackage{lmodern}
\usepackage{amsmath,amssymb,mathtools}
\usepackage{booktabs}
\usepackage[table]{xcolor}
\definecolor{SurfaceAmber}{HTML}{FFD166}
\definecolor{SurfaceRed}{HTML}{C00000}
\definecolor{SurfaceBlue}{HTML}{2A78D6}
\newcommand{\heatcell}[3][black]{\cellcolor{#2}\textcolor{#1}{#3}}
\usepackage{enumitem}
\usepackage{array,tabularx}
\usepackage{titlesec}
\usepackage{needspace}
\usepackage[skip=4pt plus 1pt minus 1pt,indent=0pt]{parskip}
\usepackage[margin=1in]{geometry}
\usepackage{microtype}
\usepackage[hidelinks]{hyperref}

% Consistent hierarchy and spacing for prose and worked calculations.
\titleformat{\section}{\large\bfseries}{\thesection}{0.8em}{}
\titleformat{\subsection}{\normalsize\bfseries}{\thesubsection}{0.8em}{}
\titleformat{\subsubsection}{\normalsize\bfseries}{}{0pt}{}
\titlespacing*{\section}{0pt}{20pt plus 3pt minus 2pt}{8pt}
\titlespacing*{\subsection}{0pt}{14pt plus 2pt minus 2pt}{6pt}
\titlespacing*{\subsubsection}{0pt}{10pt plus 2pt minus 1pt}{4pt}
\setlist[itemize]{leftmargin=1.5em,topsep=4pt,itemsep=3pt,parsep=0pt}
\setlist[description]{style=nextline,leftmargin=0pt,labelindent=0pt,labelsep=0pt,
  font=\normalfont\bfseries,topsep=4pt,itemsep=7pt,parsep=2pt}
\AtBeginDocument{%
  \setlength{\abovedisplayskip}{8pt plus 2pt minus 2pt}%
  \setlength{\belowdisplayskip}{8pt plus 2pt minus 2pt}%
  \setlength{\abovedisplayshortskip}{4pt plus 2pt}%
  \setlength{\belowdisplayshortskip}{6pt plus 2pt minus 2pt}}
\renewcommand{\arraystretch}{1.15}
\setlength{\tabcolsep}{7pt}
\setlength{\jot}{5pt}
\clubpenalty=3000
\widowpenalty=3000
\displaywidowpenalty=3000
\raggedbottom
\newcommand{\topic}[1]{\Needspace{6\baselineskip}\subsubsection*{#1}}
\newenvironment{surfaceexample}{%
  \par\Needspace{12\baselineskip}\begingroup
  \renewcommand{\arraystretch}{1.25}\setlength{\arraycolsep}{6pt}%
}{\endgroup\par}

% Highlight only conceptual anchors; body text retains its normal scale.
\definecolor{ConceptInk}{HTML}{C2410C}
\colorlet{ConceptPaper}{ConceptInk!6!white}
\newsavebox{\conceptbox}
\newenvironment{keyconcept}[1]{%
  \par\addvspace{8pt}\begin{lrbox}{\conceptbox}%
  \begin{minipage}{\dimexpr\linewidth-24pt\relax}%
  \setlength{\parindent}{0pt}%
  {\normalsize\bfseries\color{ConceptInk}#1\par}\large
}{%
  \end{minipage}\end{lrbox}%
  \noindent{\setlength{\fboxsep}{12pt}%
  \colorbox{ConceptPaper}{\usebox{\conceptbox}}}\par\addvspace{8pt}}

\newcommand{\ind}{\mathbf{1}}
\DeclareMathOperator{\Close}{Close}
\DeclareMathOperator{\Open}{Open}
\DeclareMathOperator{\High}{High}
\DeclareMathOperator{\Low}{Low}
\DeclareMathOperator{\chol}{chol}
\DeclareMathOperator{\quantile}{quantile}
\DeclareMathOperator{\searchsorted}{searchsorted}

\title{BarrierLab: Mathematical Specification}
\author{}
\date{Measurement version \texttt{shared-outcome-cache-float64-v2}}

\begin{document}
\maketitle
\tableofcontents
\clearpage

\section{Purpose and scope}

BarrierLab measures how conditioning on a feature changes the probability
of touching signed return barriers over forward horizons. Its basic
computational unit is a complete barrier--horizon Surface. The four stages
measure probability surfaces, subtract a market baseline, validate a score
of each resulting shift surface against synthetic histories, and retain the
observed surfaces that clear the test.

This document defines those transformations and their numerical contract.
The same measurement and scoring rules apply to observed and synthetic
histories. Section~\ref{sec:numerics} collects array layouts, artifact
references, and cache rules; these determine how results are represented
and reused without changing the mathematical operations.

\section{Central concepts}\label{sec:surface}

\subsection{The Surface as the computational unit}

Fix an ordered signed-barrier grid $\Delta$ and an ordered forward-horizon
grid $\Theta$. Their elements are the two coordinates of a Surface:
\begin{itemize}
 \item $\delta\in\Delta$ is a signed return barrier, expressed as a
       fraction of the starting date's close. For example, $\delta=-0.05$
       is a barrier $5\%$ below that close, and $\delta=+0.05$ is a barrier
       $5\%$ above it.
 \item $\theta\in\Theta$ is a forward horizon measured in subsequent
       observations (bars). For example, $\theta=10$ asks whether the
       barrier is touched within the next ten bars, rather than only
       at the end of the tenth bar.
\end{itemize}
Thus $\Delta$ and $\Theta$ are the grids, while $\delta$ and $\theta$
select one barrier and one horizon within them. A real-valued Surface
is a function on their product:
\begin{keyconcept}{The Surface space}
\[
  \mathcal S=\mathbb R^{\Delta\times\Theta},
  \qquad F\in\mathcal S,
  \qquad F(\delta,\theta)\in\mathbb R.
\]
\end{keyconcept}
Thus $\mathcal S$ is the space of surfaces, and $F$ is one Surface.
Its array representation has $|\Delta|$ rows and $|\Theta|$ columns.
The entire grid is one element---one point in this space---even though its
construction and evaluation use individual coordinates.

\topic{Paths and outer indices}

The path index $i$ identifies one complete OHLCV price history.
Path $i=0$ is the observed market history. Paths $i=1,\ldots,I$ are
random synthetic histories generated from the fitted null model for
hypothesis testing; the current configuration uses $I=1{,}000$.
Each path is measured and scored by the same procedure, so its surfaces
and scores can be compared with those of the observed history.
Section~\ref{sec:null} defines how these null histories are generated.

We write the path label as a parenthesized superscript $(i)$, and the
within-path date or bin as a subscript. The parentheses distinguish the
path label from an exponent. For example, $X_{t}^{(i)}$ is
the outcome Surface for path $i$ and date $t$, whereas $Y_{k}^{(i)}$ is the
conditional probability Surface for path $i$ and bin $k$. A date-level
outcome and a bin-level probability do not have the same outer indices:
aggregation replaces a collection indexed by dates with one indexed by bins.

\topic{Partial support}

A complete Surface retains the full grid even when some entries are
unsupported. Formally, a partially supported Surface is represented by
an element $F\in\mathcal S$ together with a Boolean support mask on
$\Delta\times\Theta$. Values outside the support are arbitrary and never
used. Numerical arrays may store non-finite values there instead;
NaN is a storage convention, not an element of $\mathbb R$. Completeness
describes the grid, not a guarantee that every entry is available.

\subsection{The objects carried through the pipeline}

All four Surface objects belong to the same ambient space:
\begin{keyconcept}{Four Surface objects, one common space}
\[
  X_{t}^{(i)},\quad Y_{k}^{(i)},\quad B^{(i)},\quad Z_{k}^{(i)}\in\mathcal S,
\]
\end{keyconcept}
with the support-mask convention above wherever needed. Their outer
indices identify different members of $\mathcal S$; their common internal
coordinates are always $(\delta,\theta)$. On supported coordinates:
\begin{description}[font=\normalfont\large\bfseries\color{ConceptInk}]
 \item[Outcome Surface] $X_{t}^{(i)}(\delta,\theta)\in\{0,1\}$ records whether the signed
       barrier $\delta$ was touched within horizon $\theta$ after date $t$:
       $0$ means not touched and $1$ means touched.
 \item[Conditional probability Surface] $Y_{k}^{(i)}(\delta,\theta)\in[0,1]$ gives the measured probability
       of that touch among eligible dates in condition bin $k$.
       The interval $[0,1]$ corresponds to
       $0\%$--$100\%$; for example, $0.60$ means a $60\%$ touch probability.
 \item[Unconditional probability Surface (baseline Surface)] $B^{(i)}(\delta,\theta)\in[0,1]$ gives the unconditional
       probability of the same touch, using every price-eligible date
       in path $i$, regardless of feature values or bin membership.
       It uses the same probability scale: $0.40$ means $40\%$.
       There is one unconditional probability Surface per path, shared by all condition bins.
 \item[Shift Surface] $Z_{k}^{(i)}(\delta,\theta)= Y_{k}^{(i)}(\delta,\theta)-B^{(i)}(\delta,\theta)$
       is the change in touch probability, as a probability difference.
       Its range is $[-1,1]$: positive means more likely to touch,
       negative means less likely, and zero means the same probability
       as the baseline. For example, $0.60$ conditional versus $0.40$
       baseline gives $+0.20$, displayed as $+20\%$: a 20-percentage-point
       difference, not a relative percentage change.
\end{description}
These different value ranges and meanings do not change the ambient space.

The baseline is the unrestricted case of the same averaging operation
that produces conditional probability surfaces. Writing $\Pi^{(i)}[A]$ for
the probability Surface measured over a date group $A$, we have
\[
 Y_{k}^{(i)}=\Pi^{(i)}[A_{k}^{(i)}],\qquad B^{(i)}=\Pi^{(i)}[\mathcal T],
\]
where $A_{k}^{(i)}$ contains the dates assigned to bin $k$, and $\mathcal T$
contains all observation dates. Horizon-specific price eligibility applies
in both cases. Section~\ref{sec:averaging} defines this operator precisely.
The baseline is not bin $0$: that position is already the lowest feature
bin and requires a finite feature value.

\begin{center}
\large
\renewcommand{\arraystretch}{1.45}
\begin{tabularx}{\linewidth}{@{}l l >{\raggedright\arraybackslash}X@{}}
\toprule
\rowcolor{ConceptPaper}
\textbf{Object} & \textbf{One Surface per} & \textbf{Meaning of an entry} $(\delta,\theta)$ \\
\midrule
$X_{t}^{(i)}$ & Path and date & Whether the barrier was touched \\
$Y_{k}^{(i)}$ & Path and bin & Conditional touch probability \\
$B^{(i)}$ & Path & Unconditional touch probability \\
$Z_{k}^{(i)}$ & Path and bin & Conditional minus baseline \\
\bottomrule
\end{tabularx}
\end{center}

\subsection{The scalar score and its validity flag}\label{sec:score-objects}

Scoring assigns two scalar quantities to each path $i$ and condition bin $k$:
\begin{keyconcept}{One scalar score and one validity flag per bin}
\[
 Q_k^{(i)}\in\mathbb R_{\geq0}\subset\mathbb R,
 \qquad V_k^{(i)}\in\{0,1\}.
\]
\end{keyconcept}
Their outer indices identify a path and bin, just as for $Z_k^{(i)}$.
They have no internal barrier or horizon coordinates: each summarizes
one complete shift Surface.

\begin{description}[font=\normalfont\large\bfseries\color{ConceptInk}]
 \item[Score] $Q_k^{(i)}$ is a nonnegative real number obtained by summing
       weighted contributions from usable coordinates of $Z_k^{(i)}$.
       It measures the magnitude of the difference between shifts at
       opposite signed barriers. It is not a probability and can exceed one.
 \item[Validity flag] $V_k^{(i)}$ is a Boolean indicator encoded as $0$ or $1$.
       It equals $1$ when the score has at least one usable coordinate,
       and $0$ when it has none. An unsupported score is stored as
       $Q_k^{(i)}=0$ with $V_k^{(i)}=0$; a supported score may also be zero,
       with $V_k^{(i)}=1$.
\end{description}
Validity records support, not statistical significance. Section~\ref{sec:score}
defines the contributions and support rule; Stage~3 compares supported
observed scores with scores from synthetic histories.

\begin{keyconcept}{The central computation}
\normalsize
\[
 \{X_{t}^{(i)}\}_{t\in\mathcal T}
 \xrightarrow{\text{eligible averaging}}
 \bigl(\{Y_{k}^{(i)}\}_k,B^{(i)}\bigr)
 \xrightarrow{\text{baseline subtraction}}
 \{Z_{k}^{(i)}\}_k
 \xrightarrow{\text{scoring}}
 \{(Q_{k}^{(i)},V_{k}^{(i)})\}_k.
\]
\end{keyconcept}
Averaging uses horizon-dependent eligible dates, so different
columns need not have the same denominator.

Stage~1 constructs the observed probability surfaces. Stage~2 constructs
their shifts $Z_{k}^{(0)}$. Stage~3 scores the observed shifts and repeats the
computation on synthetic histories to obtain a null comparison for each
bin. Stage~4 selects and renders the complete observed shift surfaces
whose bins clear that comparison.

The Surface is the unit being scored and selected. Its score also depends
on support in the other bins of the same feature: a barrier pair contributes
only when at least two bins are usable at that horizon. This dependency is
made explicit in Section~\ref{sec:score}.

\section{Mathematical objects and conventions}\label{sec:notation}

\subsection{Outer indices and Surface coordinates}

The exposition fixes one feature (one condition node). The computation is
repeated for every node; feature-dependent quantities below implicitly
belong to that fixed node. Outcome surfaces and market baselines are
shared across nodes within a path.

The path label always appears as a parenthesized superscript; subscripts
identify dates, bins, or grid coordinates within that path. When a quantity
also has a descriptive superscript, both labels are retained, for example
$N_{k,\theta}^{\mathrm{condition},(i)}$ for conditional observation counts
and $R_{t,\theta}^{-,(i)}$ for downside excursions.

\begin{tabularx}{\linewidth}{@{}l >{\raggedright\arraybackslash}X@{}}
\toprule
Symbol & Meaning \\
\midrule
$i=0,1,\ldots,I$ & identifies a price-history path:
        $i=0$ is observed, and $i=1,\ldots,I$ are null replicates. \\
$t\in\mathcal T$ & identifies an observation date. In forward
        expressions, $t+j$ means the $j$th subsequent observation. \\
$x_{t}^{(i)}$ & is the feature value used to assign date $t$ to a bin. \\
$k\in\{0,\ldots,K_{\mathrm{req}}-1\}$ & is a requested
        bin position. Path $i$ has $K^{(i)}\leq K_{\mathrm{req}}$ effective
        bins, occupying positions $0,\ldots,K^{(i)}-1$. \\
\bottomrule
\end{tabularx}

Bin positions are zero-based, matching the insertion-index convention
used in the implementation. Bin $k$ refers to an ordered position within
each path's binning; its numerical feature boundaries can differ between
paths.

\Needspace{13\baselineskip}
\subsection{Values, counts, and support}

Outcome entries are Boolean on price-eligible coordinates; probability
entries lie in $[0,1]$ where supported; shift entries are probability differences in $[-1,1]$. The support threshold is
\[
  n_{\min}=30.
\]
Counts depend on horizon because fewer dates have longer forward windows.
Probability and shift surfaces may therefore be only partially supported.

A support mask is separate from a value. In particular, an unsupported
score is stored as zero with a false flag; a supported score can also equal
zero. All sums below use eligible or usable index sets so that excluded
non-finite values are never interpreted as products such as $0\cdot
\mathrm{NaN}$.

\section{Stage 1: measurement of probability surfaces}\label{sec:measurement}

Measurement turns each history into date-level outcome surfaces, then
averages those outcomes over all eligible dates and over feature bins.
Its outputs are conditional probability surfaces $Y_{k}^{(i)}$,
an unconditional probability Surface $B^{(i)}$, and
the counts supporting them. Stage~1 retains the observed outputs ($i=0$);
Stage~3 applies the same definitions to null paths.

\subsection{From prices to an outcome Surface}

For every path, starting date, and horizon, define the greatest downside
and upside excursion reached by future intrabar prices:
\begin{align}
R_{t,\theta}^{-,(i)}
  &=\frac{\min_{1\leq j\leq\theta}\Low_{t+j}^{(i)}}{\Close_{t}^{(i)}}-1,\\
R_{t,\theta}^{+,(i)}
  &=\frac{\max_{1\leq j\leq\theta}\High_{t+j}^{(i)}}{\Close_{t}^{(i)}}-1.
\end{align}
Dates without $\theta$ future observations have non-finite excursions.
Define the price-eligible date set
\[
 \mathcal E_{\theta}^{(i)}
 =\{t\in\mathcal T:R_{t,\theta}^{-,(i)}
                    \text{ and }R_{t,\theta}^{+,(i)}\text{ are finite}\}.
\]
For eligible dates, the outcome Surface has entries
\[
X_{t}^{(i)}(\delta,\theta)=
\begin{cases}
 \ind[R_{t,\theta}^{-,(i)}\leq\delta],&\delta<0,\\
 \ind[R_{t,\theta}^{+,(i)}\geq\delta],&\delta\geq0.
\end{cases}
\]
An entry answers: starting from this date's close, was the signed barrier
reached by a future low or high within this horizon? The zero barrier uses
the upside branch. It belongs to the measured grid but is excluded from
the paired-barrier score.

$X_{t}^{(i)}$ depends only on the history and grids. Feature bins select which
outcome surfaces to average; they do not change the outcome surfaces.

\subsection{From feature values to date groups}

For a continuous feature, request the interior quantile edges
\[
 e_{j}^{(i)}=\quantile_{j/K_{\mathrm{req}}}
       \bigl(\{x_{t}^{(i)}:x_{t}^{(i)}\text{ is finite}\}\bigr),
 \qquad j=1,\ldots,K_{\mathrm{req}}-1.
\]
Remove non-finite and duplicate edges, and retain only edges satisfying
$\min(x^{(i)})<e_{j}^{(i)}\leq\max(x^{(i)})$, with extrema taken over finite
feature values. Write $e^{(i)}$ for the resulting ordered edge vector and
$K^{(i)}=|e^{(i)}|+1$. Heavy ties can reduce the effective bin count; an edge at
the maximum can leave the final bin empty.

For finite feature values, assign
\[
 \kappa_{t}^{(i)}=\searchsorted(e^{(i)},x_{t}^{(i)},\mathrm{left}).
\]
An observation equal to an edge enters the lower bin. Define the date group
for bin $k$ and its price-eligible subset at horizon $\theta$ by
\[
 A_{k}^{(i)}=\{t\in\mathcal T:x_{t}^{(i)}\text{ is finite},\ \kappa_{t}^{(i)}=k\},
 \qquad
 \mathcal E_{k,\theta}^{(i)}=A_{k}^{(i)}\cap\mathcal E_{\theta}^{(i)}.
\]
Positions $k\geq K^{(i)}$ have no eligible dates.

\subsection{Averaging outcome surfaces}\label{sec:averaging}

For any date group $A\subseteq\mathcal T$, define the probability Surface
\[
 \Pi^{(i)}[A](\delta,\theta)
 =\frac{\displaystyle\sum_{t\in A\cap\mathcal E_{\theta}^{(i)}}
                       X_{t}^{(i)}(\delta,\theta)}
        {|A\cap\mathcal E_{\theta}^{(i)}|},
 \qquad |A\cap\mathcal E_{\theta}^{(i)}|\geq n_{\min}.
\]
Entries are non-finite when the count requirement fails. Conditional
measurement and baseline measurement are two applications of this operator:
\[
 Y_{k}^{(i)}=\Pi^{(i)}[A_{k}^{(i)}],\qquad B^{(i)}=\Pi^{(i)}[\mathcal T].
\]
Thus the baseline is the special case with no restriction on the date
group; it still enforces price eligibility and the minimum count.

To make the supporting counts explicit, define
\begin{align}
 N_{k,\theta}^{\mathrm{condition},(i)}&=|\mathcal E_{k,\theta}^{(i)}|,\\
 N_{k}^{\mathrm{hit},(i)}(\delta,\theta)
   &=\sum_{t\in\mathcal E_{k,\theta}^{(i)}}X_{t}^{(i)}(\delta,\theta),\\
 N_{\theta}^{\mathrm{baseline},(i)}&=|\mathcal E_{\theta}^{(i)}|.
\end{align}
The probability surfaces are
\begin{align}
 Y_{k}^{(i)}(\delta,\theta)
   &=\frac{N_{k}^{\mathrm{hit},(i)}(\delta,\theta)}
           {N_{k,\theta}^{\mathrm{condition},(i)}},
   &&N_{k,\theta}^{\mathrm{condition},(i)}\geq n_{\min},\\
 B^{(i)}(\delta,\theta)
   &=\frac{\sum_{t\in\mathcal E_{\theta}^{(i)}}X_{t}^{(i)}(\delta,\theta)}
           {N_{\theta}^{\mathrm{baseline},(i)}},
   &&N_{\theta}^{\mathrm{baseline},(i)}\geq n_{\min}.
\end{align}
When the stated count requirement fails, the corresponding entries are
non-finite.

The baseline includes every price-eligible date, independently of feature
warm-up or missing feature values. The conditional surface uses only
eligible dates in its bin. Thus the baseline generally cannot be recovered
by averaging the conditional surfaces over bins. Each synthetic path has
its own baseline $B^{(i)}$.

\Needspace{20\baselineskip}
\subsection{Worked example: measuring a Surface}\label{sec:example}

Take five barriers and seven daily horizons:
\[
 \Delta=\{-2\%,-1\%,0\%,+1\%,+2\%\},
 \qquad \Theta=\{1,2,3,4,5,6,7\}.
\]
The percentage labels stand for returns $-0.02,-0.01,0,0.01,0.02$.
Assume a market with a bar every calendar day, so seven bars cover one
week. For a trading-day-only history, these horizons instead cover seven
trading days.

Start from a close of $100$. The five barrier prices are then
$98,99,100,101,102$. Suppose future intrabar prices first reach
$100$ on day~1, $101$ on day~2, $99$ on day~3, $102$ on day~5,
and $98$ on day~6. This gives the following outcome Surface:
\begin{surfaceexample}
\[
X_t^{(i)}=
\begin{array}{c|rrrrrrr}
 \delta\backslash\theta &1&2&3&4&5&6&7\\ \hline
 -2\% &\heatcell[black]{SurfaceAmber!0!white}{0}&\heatcell[black]{SurfaceAmber!0!white}{0}&\heatcell[black]{SurfaceAmber!0!white}{0}&\heatcell[black]{SurfaceAmber!0!white}{0}&\heatcell[black]{SurfaceAmber!0!white}{0}&\heatcell[white]{SurfaceRed!100!SurfaceAmber}{1}&\heatcell[white]{SurfaceRed!100!SurfaceAmber}{1}\\
 -1\% &\heatcell[black]{SurfaceAmber!0!white}{0}&\heatcell[black]{SurfaceAmber!0!white}{0}&\heatcell[white]{SurfaceRed!100!SurfaceAmber}{1}&\heatcell[white]{SurfaceRed!100!SurfaceAmber}{1}&\heatcell[white]{SurfaceRed!100!SurfaceAmber}{1}&\heatcell[white]{SurfaceRed!100!SurfaceAmber}{1}&\heatcell[white]{SurfaceRed!100!SurfaceAmber}{1}\\
  0\% &\heatcell[white]{SurfaceRed!100!SurfaceAmber}{1}&\heatcell[white]{SurfaceRed!100!SurfaceAmber}{1}&\heatcell[white]{SurfaceRed!100!SurfaceAmber}{1}&\heatcell[white]{SurfaceRed!100!SurfaceAmber}{1}&\heatcell[white]{SurfaceRed!100!SurfaceAmber}{1}&\heatcell[white]{SurfaceRed!100!SurfaceAmber}{1}&\heatcell[white]{SurfaceRed!100!SurfaceAmber}{1}\\
 +1\% &\heatcell[black]{SurfaceAmber!0!white}{0}&\heatcell[white]{SurfaceRed!100!SurfaceAmber}{1}&\heatcell[white]{SurfaceRed!100!SurfaceAmber}{1}&\heatcell[white]{SurfaceRed!100!SurfaceAmber}{1}&\heatcell[white]{SurfaceRed!100!SurfaceAmber}{1}&\heatcell[white]{SurfaceRed!100!SurfaceAmber}{1}&\heatcell[white]{SurfaceRed!100!SurfaceAmber}{1}\\
 +2\% &\heatcell[black]{SurfaceAmber!0!white}{0}&\heatcell[black]{SurfaceAmber!0!white}{0}&\heatcell[black]{SurfaceAmber!0!white}{0}&\heatcell[black]{SurfaceAmber!0!white}{0}&\heatcell[white]{SurfaceRed!100!SurfaceAmber}{1}&\heatcell[white]{SurfaceRed!100!SurfaceAmber}{1}&\heatcell[white]{SurfaceRed!100!SurfaceAmber}{1}
\end{array}.
\]
\end{surfaceexample}
For example, $X_t^{(i)}(+2\%,4)=0$ but $X_t^{(i)}(+2\%,5)=1$:
the price had not reached $102$ within four days, but had within five.
Once a row becomes $1$, it stays $1$ at longer horizons, even if price
later reverses. The zero row records a future high at or above the starting
close; it is not automatically a touch at the starting date.

Now collect 100 starting dates assigned to the same condition bin $k$,
each with seven eligible future days. Each date contributes one such
$5\times7$ outcome Surface. Suppose their coordinatewise averages are
\begin{surfaceexample}
\[
Y_k^{(i)}=
\begin{array}{c|rrrrrrr}
 \delta\backslash\theta &1&2&3&4&5&6&7\\ \hline
 -2\% &\heatcell[black]{SurfaceAmber!10!white}{0.05}&\heatcell[black]{SurfaceAmber!20!white}{0.10}&\heatcell[black]{SurfaceAmber!30!white}{0.15}&\heatcell[black]{SurfaceAmber!40!white}{0.20}&\heatcell[black]{SurfaceAmber!50!white}{0.25}&\heatcell[black]{SurfaceAmber!60!white}{0.30}&\heatcell[black]{SurfaceAmber!70!white}{0.35}\\
 -1\% &\heatcell[black]{SurfaceAmber!30!white}{0.15}&\heatcell[black]{SurfaceAmber!40!white}{0.20}&\heatcell[black]{SurfaceAmber!50!white}{0.25}&\heatcell[black]{SurfaceAmber!60!white}{0.30}&\heatcell[black]{SurfaceAmber!70!white}{0.35}&\heatcell[black]{SurfaceAmber!80!white}{0.40}&\heatcell[black]{SurfaceAmber!90!white}{0.45}\\
  0\% &\heatcell[black]{SurfaceRed!40!SurfaceAmber}{0.70}&\heatcell[black]{SurfaceRed!50!SurfaceAmber}{0.75}&\heatcell[white]{SurfaceRed!60!SurfaceAmber}{0.80}&\heatcell[white]{SurfaceRed!70!SurfaceAmber}{0.85}&\heatcell[white]{SurfaceRed!80!SurfaceAmber}{0.90}&\heatcell[white]{SurfaceRed!90!SurfaceAmber}{0.95}&\heatcell[white]{SurfaceRed!100!SurfaceAmber}{1.00}\\
 +1\% &\heatcell[black]{SurfaceAmber!50!white}{0.25}&\heatcell[black]{SurfaceAmber!60!white}{0.30}&\heatcell[black]{SurfaceAmber!70!white}{0.35}&\heatcell[black]{SurfaceAmber!80!white}{0.40}&\heatcell[black]{SurfaceAmber!90!white}{0.45}&\heatcell[black]{SurfaceAmber!100!white}{0.50}&\heatcell[black]{SurfaceRed!10.000000000000014!SurfaceAmber}{0.55}\\
 +2\% &\heatcell[black]{SurfaceAmber!30!white}{0.15}&\heatcell[black]{SurfaceAmber!40!white}{0.20}&\heatcell[black]{SurfaceAmber!50!white}{0.25}&\heatcell[black]{SurfaceAmber!60!white}{0.30}&\heatcell[black]{SurfaceAmber!70!white}{0.35}&\heatcell[black]{SurfaceAmber!80!white}{0.40}&\heatcell[black]{SurfaceAmber!90!white}{0.45}
\end{array}.
\]
\end{surfaceexample}
The entry $Y_k^{(i)}(+2\%,5)=0.35$ means that 35 of the 100 dates
were followed by a touch of the $+2\%$ barrier within five days.
The hit count is 35, the observation count is 100, and the measured
probability is $35/100=35\%$. All 35 entries together form one
conditional probability Surface. Real histories can have different
observation counts at different horizons.

\section{Stage 2: comparison with the baseline}\label{sec:comparison}

Comparison converts a conditional probability Surface into a
baseline-relative shift Surface.

\subsection{Baseline subtraction}

Define the shift Surface by
\[
 Z_{k}^{(i)}= (Y_{k}^{(i)}-B^{(i)}).
\]
Subtraction is coordinatewise on their common support:
\[
 Z_{k}^{(i)}(\delta,\theta)
 =\bigl(Y_{k}^{(i)}(\delta,\theta)-B^{(i)}(\delta,\theta)\bigr).
\]
A positive entry means that the barrier was touched more often in the
condition bin than in the path's unconditional market history. A negative
entry means it was touched less often. Null histories use their own $B^{(i)}$,
never the observed baseline $B^{(0)}$.

\subsection{Worked example: subtracting the baseline}

Continue with the same five barriers, seven daily horizons, and conditional
Surface $Y_k^{(i)}$ from Section~\ref{sec:example}. Suppose averaging over
all price-eligible dates in the same path gives the baseline
\begin{surfaceexample}
\[
B^{(i)}=
\begin{array}{c|rrrrrrr}
 \delta\backslash\theta &1&2&3&4&5&6&7\\ \hline
 -2\% &\heatcell[black]{SurfaceAmber!30!white}{0.15}&\heatcell[black]{SurfaceAmber!40!white}{0.20}&\heatcell[black]{SurfaceAmber!50!white}{0.25}&\heatcell[black]{SurfaceAmber!60!white}{0.30}&\heatcell[black]{SurfaceAmber!70!white}{0.35}&\heatcell[black]{SurfaceAmber!80!white}{0.40}&\heatcell[black]{SurfaceAmber!90!white}{0.45}\\
 -1\% &\heatcell[black]{SurfaceAmber!50!white}{0.25}&\heatcell[black]{SurfaceAmber!60!white}{0.30}&\heatcell[black]{SurfaceAmber!70!white}{0.35}&\heatcell[black]{SurfaceAmber!80!white}{0.40}&\heatcell[black]{SurfaceAmber!90!white}{0.45}&\heatcell[black]{SurfaceAmber!100!white}{0.50}&\heatcell[black]{SurfaceRed!10.000000000000014!SurfaceAmber}{0.55}\\
  0\% &\heatcell[black]{SurfaceRed!40!SurfaceAmber}{0.70}&\heatcell[black]{SurfaceRed!50!SurfaceAmber}{0.75}&\heatcell[white]{SurfaceRed!60!SurfaceAmber}{0.80}&\heatcell[white]{SurfaceRed!70!SurfaceAmber}{0.85}&\heatcell[white]{SurfaceRed!80!SurfaceAmber}{0.90}&\heatcell[white]{SurfaceRed!90!SurfaceAmber}{0.95}&\heatcell[white]{SurfaceRed!100!SurfaceAmber}{1.00}\\
 +1\% &\heatcell[black]{SurfaceAmber!30!white}{0.15}&\heatcell[black]{SurfaceAmber!40!white}{0.20}&\heatcell[black]{SurfaceAmber!50!white}{0.25}&\heatcell[black]{SurfaceAmber!60!white}{0.30}&\heatcell[black]{SurfaceAmber!70!white}{0.35}&\heatcell[black]{SurfaceAmber!80!white}{0.40}&\heatcell[black]{SurfaceAmber!90!white}{0.45}\\
 +2\% &\heatcell[black]{SurfaceAmber!10!white}{0.05}&\heatcell[black]{SurfaceAmber!20!white}{0.10}&\heatcell[black]{SurfaceAmber!30!white}{0.15}&\heatcell[black]{SurfaceAmber!40!white}{0.20}&\heatcell[black]{SurfaceAmber!50!white}{0.25}&\heatcell[black]{SurfaceAmber!60!white}{0.30}&\heatcell[black]{SurfaceAmber!70!white}{0.35}
\end{array}.
\]
\end{surfaceexample}
Subtract each baseline entry from the conditional entry at the same
barrier and horizon. The result is one complete
shift Surface, with every entry below on the probability-difference scale:
\begin{surfaceexample}
\[
Z_k^{(i)}= (Y_k^{(i)}-B^{(i)})=
\begin{array}{c|rrrrrrr}
 \delta\backslash\theta &1&2&3&4&5&6&7\\ \hline
 -2\% &\heatcell[black]{SurfaceBlue!33.33333333333333!white}{-0.10}&\heatcell[black]{SurfaceBlue!33.33333333333333!white}{-0.10}&\heatcell[black]{SurfaceBlue!33.33333333333333!white}{-0.10}&\heatcell[black]{SurfaceBlue!33.33333333333333!white}{-0.10}&\heatcell[black]{SurfaceBlue!33.33333333333333!white}{-0.10}&\heatcell[black]{SurfaceBlue!33.33333333333333!white}{-0.10}&\heatcell[black]{SurfaceBlue!33.33333333333333!white}{-0.10}\\
 -1\% &\heatcell[black]{SurfaceBlue!33.33333333333333!white}{-0.10}&\heatcell[black]{SurfaceBlue!33.33333333333333!white}{-0.10}&\heatcell[black]{SurfaceBlue!33.33333333333333!white}{-0.10}&\heatcell[black]{SurfaceBlue!33.33333333333333!white}{-0.10}&\heatcell[black]{SurfaceBlue!33.33333333333333!white}{-0.10}&\heatcell[black]{SurfaceBlue!33.33333333333333!white}{-0.10}&\heatcell[black]{SurfaceBlue!33.33333333333333!white}{-0.10}\\
  0\% &\heatcell[black]{SurfaceRed!0!white}{0}&\heatcell[black]{SurfaceRed!0!white}{0}&\heatcell[black]{SurfaceRed!0!white}{0}&\heatcell[black]{SurfaceRed!0!white}{0}&\heatcell[black]{SurfaceRed!0!white}{0}&\heatcell[black]{SurfaceRed!0!white}{0}&\heatcell[black]{SurfaceRed!0!white}{0}\\
 +1\% &\heatcell[black]{SurfaceRed!33.33333333333333!white}{+0.10}&\heatcell[black]{SurfaceRed!33.33333333333333!white}{+0.10}&\heatcell[black]{SurfaceRed!33.33333333333333!white}{+0.10}&\heatcell[black]{SurfaceRed!33.33333333333333!white}{+0.10}&\heatcell[black]{SurfaceRed!33.33333333333333!white}{+0.10}&\heatcell[black]{SurfaceRed!33.33333333333333!white}{+0.10}&\heatcell[black]{SurfaceRed!33.33333333333333!white}{+0.10}\\
 +2\% &\heatcell[black]{SurfaceRed!33.33333333333333!white}{+0.10}&\heatcell[black]{SurfaceRed!33.33333333333333!white}{+0.10}&\heatcell[black]{SurfaceRed!33.33333333333333!white}{+0.10}&\heatcell[black]{SurfaceRed!33.33333333333333!white}{+0.10}&\heatcell[black]{SurfaceRed!33.33333333333333!white}{+0.10}&\heatcell[black]{SurfaceRed!33.33333333333333!white}{+0.10}&\heatcell[black]{SurfaceRed!33.33333333333333!white}{+0.10}
\end{array}.
\]
\end{surfaceexample}
For example, at $(+2\%,5)$ the conditional probability is $35\%$
and the baseline is $25\%$:
\[
 Z_k^{(i)}(+2\%,5)=0.35-0.25=+0.10.
\]
At $(-2\%,5)$, the corresponding subtraction is
$0.25-0.35=-0.10$. The zero-barrier probabilities are equal,
so their shifts are zero.
The condition is associated with fewer downside touches and more upside
touches across the seven horizons. Stage~2 retains the observed $Z_{k}^{(0)}$ together
with references to its measurement artifacts.

\section{Stage 3: validation by null comparison}\label{sec:validation}

Validation asks how often the specified synthetic-history model produces
a bin score at least as large as the observed score. It first reduces each
shift Surface to a scalar, then repeats measurement and scoring on null
histories, and finally compares scores at the same bin position.

\subsection{Scoring a shift Surface}\label{sec:score}

Scoring keeps the full signed-barrier--by--horizon grid. It compares each
shift with the shift at the opposite barrier, assigns a nonnegative
contribution to each original coordinate, and sums that contribution
Surface to obtain one bin score.

\topic{Reflection and weights}

Let $\Delta_{\pm}$ contain the nonzero barriers $\delta$ for which the
grid has exactly one $\delta$ row and exactly one $-\delta$ row. Define
reflection on the original grid by
\[
 (\operatorname{Ref}F)(\delta,\theta)=
 \begin{cases}
 F(-\delta,\theta),&\delta\in\Delta_{\pm},\\
 F(\delta,\theta),&\delta\notin\Delta_{\pm}.
 \end{cases}
\]
Reflection reorders rows without changing the Surface shape. When
$\Delta_{\pm}\ne\varnothing$, assign the magnitude weights
\[
 w_\delta=
 \begin{cases}
 \displaystyle\frac{|\delta|}{\max_{\gamma\in\Delta_{\pm}}|\gamma|},
      &\delta\in\Delta_{\pm},\\
 0,&\delta\notin\Delta_{\pm}.
 \end{cases}
\]
If $\Delta_{\pm}$ is empty, all weights are zero.

\topic{Usable coordinates}

Support is determined on this same grid. First identify locally eligible
bins at each nonzero matched barrier and horizon:
\[
 \mathcal L^{(i)}(\delta,\theta)=
 \left\{k\in\{0,\ldots,K_{\mathrm{req}}-1\}:
 \begin{array}{l}
 Z_k^{(i)}(\delta,\theta),Z_k^{(i)}(-\delta,\theta)\text{ are finite},\\
 N_{k,\theta}^{\mathrm{condition},(i)}\geq n_{\min}
 \end{array}\right\},
 \qquad \delta\in\Delta_{\pm}.
\]
The usable coordinates for bin $k$ are
\[
 \mathcal U_k^{(i)}
 =\{(\delta,\theta)\in\Delta_{\pm}\times\Theta:
 k\in\mathcal L^{(i)}(\delta,\theta),\
 |\mathcal L^{(i)}(\delta,\theta)|\geq2\}.
\]
Thus support still depends on the other bins of the same feature.
\topic{Contribution and score}

Define the full contribution Surface $H_k^{(i)}\in\mathcal S$ by
\[
 H_k^{(i)}(\delta,\theta)=
 \begin{cases}
 \displaystyle\frac{w_\delta}{2}
 \left|Z_k^{(i)}(\delta,\theta)
       -(\operatorname{Ref}Z_k^{(i)})(\delta,\theta)\right|,
       &(\delta,\theta)\in\mathcal U_k^{(i)},\\
 0,&\text{otherwise}.
 \end{cases}
\]
Both signed rows carry their own contribution. The factor $1/2$ accounts
for the same absolute difference appearing at both $\delta$ and $-\delta$,
preserving the score's scale. No reduced magnitude--horizon Surface is
constructed. Unsupported coordinates, zero barriers, and unmatched barriers
contribute zero.

The full-grid score and its validity flag are
\[
 Q_k^{(i)}=\sum_{\theta\in\Theta}\sum_{\delta\in\Delta}
                       H_k^{(i)}(\delta,\theta),
 \qquad V_k^{(i)}=\ind[\mathcal U_k^{(i)}\ne\varnothing].
\]
Equal shifts at opposite barriers contribute zero; the absolute value
discards the direction of the asymmetry. Larger barrier magnitudes receive
greater linear weight. The sum is not normalized by the number of usable
coordinates. It is a weighted probability-difference score, not a probability,
and can exceed one. An unsupported bin has score zero and false validity;
a supported zero score remains valid.

\subsection{Worked example: scoring the shift}\label{sec:score-example}

For the running example, take the Stage~2 Surface as the observed
Surface $Z_k^{(0)}$. Its two negative-barrier rows are $-0.10$ and
its two positive-barrier rows are $+0.10$ at every horizon.
The bin has 100 eligible dates at every horizon; assume at least one
other bin is locally eligible at both magnitudes and all seven horizons.
Each nonzero coordinate is therefore usable.

The $1\%$ rows have weight $0.01/0.02=0.5$, so each entry contributes
$\tfrac12(0.5)|-0.10-(+0.10)|=0.05$. The $2\%$ rows have weight $1$, so
each entry contributes $\tfrac12(1)|-0.10-(+0.10)|=0.10$. The resulting
contribution Surface has the same $5\times7$ layout as the shift Surface:
\begin{surfaceexample}
\[
H_k^{(0)}=
\begin{array}{c|rrrrrrr}
 \delta\backslash\theta &1&2&3&4&5&6&7\\ \hline
 -2\% &0.10&0.10&0.10&0.10&0.10&0.10&0.10\\
 -1\% &0.05&0.05&0.05&0.05&0.05&0.05&0.05\\
  0\% &0&0&0&0&0&0&0\\
 +1\% &0.05&0.05&0.05&0.05&0.05&0.05&0.05\\
 +2\% &0.10&0.10&0.10&0.10&0.10&0.10&0.10
\end{array}.
\]
\end{surfaceexample}
Each column contributes $0.10+0.05+0+0.05+0.10=0.30$. Summing the entire Surface gives
\[
 Q_k^{(0)}=\underbrace{0.30+0.30+0.30+0.30+0.30+0.30+0.30}_{\text{seven horizons}}
         =2.10,
 \qquad V_k^{(0)}=1.
\]
Stage~3 compares this single score with scores from null histories; it
does not test each of the 35 Surface entries separately.

\subsection{The synthetic-history model}\label{sec:null}

The null model fits the contemporaneous mean and covariance of four
observed OHLC log components, then draws independent innovations across
time. It provides the reference histories for the score comparison.

\topic{Fitting and drawing log components}

For the observed inputs below, suppress the path superscript $(0)$.
Component labels such as $\mathrm{open}$ share the superscript with the
path label when both are present, as in $g_t^{\mathrm{open},(i)}$. Define
\begin{align}
 g^{\mathrm{open}}_t
   &=\log\left(\frac{\Open_t}{\Close_{t-1}}\right),&
 g^{\mathrm{close}}_t
   &=\log\left(\frac{\Close_t}{\Open_t}\right),\\
 g^{\mathrm{high}}_t
   &=\log\left(\frac{\High_t}{\max(\Open_t,\Close_t)}\right),&
 g^{\mathrm{low}}_t
   &=\log\left(\frac{\Low_t}{\min(\Open_t,\Close_t)}\right).
\end{align}
After removing non-finite rows, let $\mu$ and $\Sigma$ be the empirical
mean vector and covariance matrix of these components. Add diagonal
jitter before Cholesky factorization:
\[
 \epsilon=10^{-12}\max(1,\max_j|\Sigma_{jj}|),
 \qquad L_\Sigma=\chol(\Sigma+\epsilon I_4).
\]
Using deterministic seed \texttt{20260907}, draw row vectors
\[
 z_{t}^{(i)}\sim\mathcal N(0,I_4),
 \qquad g_{t}^{(i)}=z_{t}^{(i)}L_\Sigma^{\mathsf T}+\mu.
\]
The draws retain the fitted contemporaneous covariance, with the stated
diagonal jitter, and are independent across time.

\topic{Reconstructing a price history}

Let $C_{\mathrm{init}}$ denote the observed close used to initialize
each synthetic history, with $\Close_{0}^{(i)}=C_{\mathrm{init}}$.
For synthetic steps $t\geq1$, reconstruct
\begin{align}
 \Close_{t}^{(i)}
   &=C_{\mathrm{init}}\exp\left(
        \sum_{j=1}^{t}(g_{j}^{\mathrm{open},(i)}+g_{j}^{\mathrm{close},(i)})
                           \right),\\
 \Open_{t}^{(i)}&=\Close_{t-1}^{(i)}\exp(g_{t}^{\mathrm{open},(i)}),\\
 \High_{t}^{(i)}&=\max(\Open_{t}^{(i)},\Close_{t}^{(i)})\exp(g_{t}^{\mathrm{high},(i)}),\\
 \Low_{t}^{(i)}&=\min(\Open_{t}^{(i)},\Close_{t}^{(i)})\exp(g_{t}^{\mathrm{low},(i)}).
\end{align}
Observed volume is broadcast across synthetic histories. The null does
not model temporal autocorrelation, volatility clustering, regime
transitions, liquidity, execution, or trading costs.

\subsection{Repeating measurement on each null path}

For an OHLCV-derived feature, recompute its values and quantile edges
independently on each synthetic path. External, calendar, and
workspace-extension conditions that cannot be reconstructed from synthetic
OHLCV reuse their observed feature values and observed edges on every
null path.

Apply Sections~\ref{sec:measurement}--\ref{sec:comparison} and
Section~\ref{sec:score} to every path. This produces
$Q_{k}^{(1)},\ldots,Q_{k}^{(I)}$ at each requested bin position. The comparison
matches bin positions; it does not require identical numerical quantile
edges or identical effective bin counts across histories.

\subsection{The test and decision rule}

For each supported observed bin, the one-sided Monte Carlo i-value is
\[
 \widehat i_k=
 \frac{1+\sum_{i=1}^{I}\ind[Q_{k}^{(i)}\geq Q_{k}^{(0)}]}{I+1}.
\]
Invalid null bins retain their zero scores in this complete ensemble.
They are not dropped from the numerator or denominator. The output
separately records $\sum_{i=1}^I V_{k}^{(i)}$, the number of valid null scores.

With the current $I=1{,}000$, the denominator is $1{,}001$.
A supported observed bin clears precisely when $\widehat i_k<0.05$.
Unsupported observed bins are skipped. There is no family-wise or
false-discovery-rate correction across bins in the current decision rule.
The comparison concerns this score under the specified null model.

\Needspace{16\baselineskip}
\subsection{Worked example: applying the decision rule}\label{sec:decision-example}

To finish the running example, the observed score is $Q_k^{(0)}=2.10$.
Generate and score 1,000 null histories using the same rules. Suppose,
for illustration, that 24 null scores are at least 2.10, including any
ties. Then
\[
 \widehat i_k=\frac{1+24}{1{,}001}
            =\frac{25}{1{,}001}\approx0.025<0.05.
\]
This bin clears, and Stage~4 retains its complete $-0.10/0/+0.10$ shift
Surface. The count of 24 is an illustrative assumption, not a simulation
result: the observed shifts alone cannot determine the null comparison.

\section{Stage 4: selection and presentation}

Stage~4 consumes the Stage~3 decision and introduces no new score.

\subsection{Selection and rendering rules}

Writing $C_k$ for the \texttt{cleared} indicator, select exactly
\[
 \mathcal K_{\mathrm{sel}}
 =\{k:C_k=\mathrm{true}\}
 =\{k:V_{k}^{(0)}=1\text{ and }\widehat i_k<0.05\}.
\]
For every selected bin, render the complete observed shift Surface
$Z_{k}^{(0)}$ as a signed-barrier--by--horizon heatmap.

There is no top-$k$ truncation, family quota, effect-size gate, or
representative-cell rule. Ordering by ascending $\widehat i_k$ and
descending $Q_{k}^{(0)}$ is presentational only.

\subsection{Worked example: selecting the shift Surface}

In the running example, the observed bin is supported and its illustrative
Monte Carlo i-value is $25/1{,}001<0.05$
(Section~\ref{sec:decision-example}). Hence $C_k=\mathrm{true}$ and
$k\in\mathcal K_{\mathrm{sel}}$.

Render the complete observed $5\times7$ shift Surface: both
negative-barrier rows contain $-0.10$, the zero-barrier row contains
$0$, and both positive-barrier rows contain $+0.10$ across all seven
horizons. The score $Q_k^{(0)}=2.10$ summarizes this Surface for validation;
the selected output retains all 35 entries. Selection of any other bin
depends on that bin's own support and decision result.

\section{Numerical and artifact contract}\label{sec:numerics}

\subsection{Surface notation and stored axes}

Surface notation groups the two internal coordinates for mathematical
reasoning; it does not prescribe contiguous axes in a stored tensor.
The correspondence with scalar-index notation is
\begin{align*}
 X_{t}^{(i)}(\delta,\theta)&=X_{t,\theta,\delta}^{(i)},\\
 Y_{k}^{(i)}(\delta,\theta)&=\pi_{\delta,k,\theta}^{\mathrm{condition},(i)},\\
 B^{(i)}(\delta,\theta)&=\pi_{\delta,\theta}^{\mathrm{baseline},(i)},\\
 Z_{k}^{(i)}(\delta,\theta)&=Z_{\delta,k,\theta}^{(i)},\\
 N_{k}^{\mathrm{hit},(i)}(\delta,\theta)&=N_{\delta,k,\theta}^{\mathrm{hit},(i)}.
\end{align*}
The implementation uses the following axis order. A path batch contains a
subset of the paths $i=0,1,\ldots,I$; an observed role batch contains one path.

\begin{center}
\begin{tabular}{ll}
\toprule
Quantity & Array axes, in order \\
\midrule
OHLCV histories & path, observation, OHLCV component \\
Excursions at one horizon & path, observation \\
Touches at one horizon & path, observation, barrier \\
Conditional probabilities and hit counts & path, barrier, bin, horizon \\
Unconditional probabilities & path, barrier, horizon \\
Shift surfaces & path, barrier, bin, horizon \\
Bin observation counts & path, bin, horizon \\
Bin scores and validity flags & path, bin \\
\bottomrule
\end{tabular}
\end{center}

OHLCV component order is open, high, low, close, volume. Feature values
have path (or one shared path) and observation axes; edges have path
(or one shared path) and edge axes; assignments have path and observation
axes. A complete logical excursion tensor adds the horizon axis.
The numerical computation can stream that axis without materializing
every complete outcome Surface simultaneously.

Batched results retain $K_{\mathrm{req}}$ bin positions, with
$k\geq K^{(i)}$ unsupported. Compact observed artifacts retain $K^{(0)}$ bins.
The default NASDAQ configuration has $|\Delta|=41$, $|\Theta|=30$,
$K_{\mathrm{req}}=10$, and $I=1{,}000$ null replicates.

\subsection{Artifacts and exact reuse}

\topic{Stored artifacts}

Stage~1 stores the observed conditional and unconditional probability surfaces,
hit counts, and bin observation counts. Observed forward excursions,
touches, and unconditional probabilities are retained as float64 source-level
cache tensors. Feature values, quantile edges, and integer bin assignments
are retained in the node artifact.

Shift arrays use \texttt{probability\_shift} in probability-difference units.
A separate shift version governs legacy conversion and invalidates old
validation results without changing Stage~1 measurements. Percentage
formatting is applied only in presentation.

Stage~2 stores only the shifts plus cryptographic references to its node
and baseline Stage~1 artifacts. Probabilities, counts, axes, conditions,
and market history are materialized from those references when required.

\topic{Shared computation and cache validity}

Horizons are traversed incrementally: running future minima and maxima
extend the state for shorter horizons. Synthetic histories are evaluated
in memory-bounded path batches. At each batch and horizon, one touch
matrix serves the baseline and all node-specific bin reductions.

The following quantities are reused exactly rather than approximated:
\begin{itemize}
 \item observed forward excursions, touches, and unconditional probabilities;
 \item observed feature values, quantile edges, and bin assignments;
 \item the synthetic touch matrix shared across the baseline and all
       nodes in a batch; and
 \item common rolling feature primitives within a history or synthetic batch.
\end{itemize}
Observed caches are keyed by ordered market-history content and validated
against the measurement version, barrier grid, and horizon grid.
Stage~2 references are protected by SHA-256 hashes. A changed history,
grid, numerical version, or referenced Stage~1 artifact invalidates reuse.

\end{document}
